\section{Discussion}
\label{sec:discussion}

\subsection{Key Findings}

\textbf{Benchmark saturation is a phase transition, not a gradual decline.} Prior work modeled saturation as a smooth monotonic trend~\cite{Liao2022,Akhtar2026}; our results demonstrate it is a discrete structural break detectable by standard change-point analysis. A threshold supports decision-making (retire benchmarks above it); a trend does not.

\textbf{The two regimes are quantitatively extreme.} Pre-breakpoint variance ($3.81\times$ global) and post-breakpoint variance ($0.20\times$ pre-segment) differ by approximately $19\times$. This is visible in the raw scatter (Figure~\ref{fig:scatter}) --- not a subtle statistical effect.

\textbf{Three independent experiments corroborate the same mechanism.} H-E1 detects the break; H-M1 confirms the pre-breakpoint character; H-M2 confirms the post-breakpoint character. Each uses pre-specified metrics and independent statistical tests.

\textbf{Interpretation of OLS reversal ($\rho$: $-0.28 \to +0.137$).} The reversal is not a contradiction --- but its cause is not established with certainty. We hypothesize it reflects PwC benchmark database composition changes between snapshots ($N=111 \to N=115$): if newly added benchmarks cluster at high paper\_count with high CoV, they would locally reverse the aggregate trend slope. This is a plausible interpretation given the snapshot sensitivity documented in L3, but it cannot be confirmed without per-snapshot benchmark identity data. PELT's residual-based approach is insulated from this reversal regardless of its origin.

\subsection{Limitations}
\label{sec:limitations}

\textbf{L1: Small pre-segment ($n_{\text{pre}}=8$) limits directional characterization.} The structural break falls near the left boundary of the paper\_count distribution. Moment-based inference (skewness, H-M3 M1/M3) is unreliable at $n=8$. Primary claims (structural break, variance compression) rely on statistics robust to small segment sizes. This is a consequence of where the natural breakpoint falls, not a methodological flaw.

\textbf{L2: Cross-sectional design cannot establish causality.} We pool benchmarks at a single time point. The two-regime structure is consistent with Goodhart saturation dynamics, but cross-sectional analysis cannot rule out alternative explanations (benchmark selection effects, metric heterogeneity across task types, benchmark age confounds). Longitudinal within-benchmark CoV trajectories would be required for causal attribution.

\textbf{L3: Dataset snapshot sensitivity.} paper\_count$^*=39$ reflects the Aug 2026 PwC snapshot. The OLS trend direction already reversed between Phase 1 ($N=111$) and our analysis ($N=115$), demonstrating snapshot sensitivity. Periodic re-analysis is recommended.

\textbf{L4: Bootstrap CI width $= 31.5$.} The 95\% CI $[38, 69.5]$ exceeds the soft $\leq 20$-paper criterion, driven by the small pre-segment. The point estimate (39) is actionable; the CI should be reported in any application.

\textbf{L5: Domain stratification not tested.} A global paper\_count$^*$ is reported. Whether task types (image classification, NLP) have different saturation thresholds is an explicit open question.

\textbf{L6: Multiple testing.} We test three primary hypotheses (H-E1, H-M1, H-M2) without a family-wise correction. Under Bonferroni ($\alpha/3 \approx 0.017$), the H-E1 permutation $p=0.035$ does not survive correction; the corroborating piecewise $F$-test ($p=0.0021$) does. We recommend treating the permutation test as one of two convergent lines of evidence (together with $p=0.0021$) rather than a standalone significance claim.

\subsection{Broader Impact}

This work provides a methodology for operationalizing benchmark retirement decisions using existing leaderboard data. The threshold paper\_count$^* \approx 39$ should be interpreted as a benchmark health indicator, not a hard retirement trigger. Benchmarks above this threshold merit closer review including rank reversal rates, standard deviation of top-$k$ methods, and qualitative assessment.

Potential misuse: mechanically retiring benchmarks above paper\_count$^*=39$ without qualitative review could prematurely sunset benchmarks where post-saturation reflects genuine consensus rather than Goodhart gaming. Retirement decisions should combine quantitative indicators with community input.
